\documentclass[10pt,a4paper]{amsart}
\usepackage{amsmath,amssymb,amsthm,graphicx}
% For commutative diagrams we use Xy-pic
\usepackage[all]{xy}

% We use 2cell for 2-commutative diagrams.
\xyoption{2cell}
% Theorem environments.
%
\theoremstyle{plain}
\newtheorem{theorem}[subsection]{Theorem}
\newtheorem{proposition}[subsection]{Proposition}
\newtheorem{lemma}[subsection]{Lemma}

\theoremstyle{definition}
\newtheorem{definition}[subsection]{Definition}
\newtheorem{example}[subsection]{Example}
\newtheorem{exercise}[subsection]{Exercise}
\newtheorem{situation}[subsection]{Situation}

\theoremstyle{remark}
\newtheorem{remark}[subsection]{Remark}
\newtheorem{remarks}[subsection]{Remarks}

\numberwithin{equation}{subsection}

% Macros
%
\def\lim{\mathop{\mathrm{lim}}\nolimits}
\def\colim{\mathop{\mathrm{colim}}\nolimits}
\def\Spec{\mathop{\mathrm{Spec}}}
\def\Hom{\mathop{\mathrm{Hom}}\nolimits}
\def\Ext{\mathop{\mathrm{Ext}}\nolimits}
\def\SheafHom{\mathop{\mathcal{H}\!\mathit{om}}\nolimits}
\def\SheafExt{\mathop{\mathcal{E}\!\mathit{xt}}\nolimits}
\def\Sch{\mathit{Sch}}
\def\Mor{\mathop{\mathrm{Mor}}\nolimits}
\def\Ob{\mathop{\mathrm{Ob}}\nolimits}
\def\Sh{\mathop{\mathit{Sh}}\nolimits}
\def\NL{\mathop{N\!L}\nolimits}
\def\CH{\mathop{\mathrm{CH}}\nolimits}
\def\proetale{{pro\text{-}\acute{e}tale}}
\def\etale{{\acute{e}tale}}
\def\QCoh{\mathit{QCoh}}
\def\Ker{\mathop{\mathrm{Ker}}}
\def\Im{\mathop{\mathrm{Im}}}
\def\Coker{\mathop{\mathrm{Coker}}}
\def\Coim{\mathop{\mathrm{Coim}}}

% Boxtimes
%
\DeclareMathSymbol{\boxtimes}{\mathbin}{AMSa}{"02}

%
% Macros for moduli stacks/spaces
%
\def\QCohstack{\mathcal{QC}\!\mathit{oh}}
\def\Cohstack{\mathcal{C}\!\mathit{oh}}
\def\Spacesstack{\mathcal{S}\!\mathit{paces}}
\def\Quotfunctor{\mathrm{Quot}}
\def\Hilbfunctor{\mathrm{Hilb}}
\def\Curvesstack{\mathcal{C}\!\mathit{urves}}
\def\Polarizedstack{\mathcal{P}\!\mathit{olarized}}
\def\Complexesstack{\mathcal{C}\!\mathit{omplexes}}
% \Pic is the operator that assigns to X its picard group, usage \Pic(X)
% \Picardstack_{X/B} denotes the Picard stack of X over B
% \Picardfunctor_{X/B} denotes the Picard functor of X over B
\def\Pic{\mathop{\mathrm{Pic}}\nolimits}
\def\Picardstack{\mathcal{P}\!\mathit{ic}}
\def\Picardfunctor{\mathrm{Pic}}
\def\Deformationcategory{\mathcal{D}\!\mathit{ef}}
\usepackage{cleveref}
\crefname{originalreference}{original reference}{original references}
\Crefname{originalreference}{Original reference}{Original references}
\makeatletter
\newlabel{lemma-base-change-does-not-hold}{{0F08}{source}{}{sourcepointer-1}{}}
\newlabel{lemma-base-change-does-not-hold@cref}{{[originalreference][0][]0F08}{[1][0][]source}{}{}{}}
\makeatother
\crefname{originalreference}{original reference}{original references}
\Crefname{originalreference}{Original reference}{Original references}
\makeatletter
\newlabel{lemma-smooth-base-change-separably-closed}{{0F0B}{source}{}{sourcepointer-2}{}}
\newlabel{lemma-smooth-base-change-separably-closed@cref}{{[originalreference][0][]0F0B}{[1][0][]source}{}{}{}}
\makeatother
\crefname{originalreference}{original reference}{original references}
\Crefname{originalreference}{Original reference}{Original references}
\makeatletter
\newlabel{remark-base-change-holds}{{0F02}{source}{}{sourcepointer-3}{}}
\newlabel{remark-base-change-holds@cref}{{[originalreference][0][]0F02}{[1][0][]source}{}{}{}}
\makeatother
\title{The first degree in which base change fails can be detected over an affine normal base with algebraically closed function field}
\author{The Stacks Project Authors}
\date{stacks-project:etale-cohomology.tex:local-snapshot-20261011; GNU FDL 1.2 or later}
\begin{document}
\maketitle
\noindent\textit{Editorial note.} All delivered mathematical text is the author's own, copied from the pinned Stacks Project chapter file etale-cohomology.tex; no statement or proof step was rewritten, reordered or removed.\par\medskip
\noindent\textit{Editorial note.} The counted proof is the author's complete argument: the proof is written out in full: the choice of the witness, the passage to the function field with the extended diagram, the identification via the invariance lemma, and the transfer of the nonvanishing through the relative Leray spectral sequence.\par\medskip
\noindent\textit{Editorial note.} The label lemma-base-change-does-not-hold-post names the counted statement, and the proof environment immediately following it in the source is delivered verbatim as its proof.\par\medskip
\noindent\textit{Editorial note.} The counted statement is carried by the author's own lemma environment in the chapter source.\par\medskip
\noindent\textit{Editorial note.} The statement's two diagrams are the author's and are delivered verbatim; the three cited inputs are kept as external pointers with their canonical Stacks tags.\par\medskip
\noindent\textit{Allowed standard prerequisites (editorial).}\begin{itemize}
\item The diagram and coefficient sheaf witnessing the failure of BC(f, n, q). Same chapter, the lemma producing a base change diagram and a coefficient sheaf for which the q-th base change map fails to be an isomorphism. Supplies the starting diagram, which the present lemma refines.
\item Invariance under extensions of separably closed fields. Same chapter, the lemma showing that cohomology with invertible torsion coefficients does not change under such an extension. Supplies the identification R(Y -> Y')\_* Z/dZ = Z/dZ used to replace the base by Spec(K).
\end{itemize}\medskip
\ifdefined\hypertarget\hypertarget{sourcepointer-1}{}\fi
\noindent\textit{Source reference.} Original reference 0F08: sources/stacks-etale-cohomology/pinned-original-source.tex, source lines 17296--17296. Prerequisite pointer invoked at proof line(s) 18124; it names a standard fact of the same chapter that the counted argument uses, and every step of the counted proof built on it is retained verbatim in the delivered body.\par\medskip
\ifdefined\hypertarget\hypertarget{sourcepointer-2}{}\fi
\noindent\textit{Source reference.} Original reference 0F0B: sources/stacks-etale-cohomology/pinned-original-source.tex, source lines 18066--18066. Prerequisite pointer invoked at proof line(s) 18141; it names a standard fact of the same chapter that the counted argument uses, and every step of the counted proof built on it is retained verbatim in the delivered body.\par\medskip
\ifdefined\hypertarget\hypertarget{sourcepointer-3}{}\fi
\noindent\textit{Source reference.} Original reference 0F02: sources/stacks-etale-cohomology/pinned-original-source.tex, source lines 16848--16848. Prerequisite pointer invoked at proof line(s) 18107; it names a standard fact of the same chapter that the counted argument uses, and every step of the counted proof built on it is retained verbatim in the delivered body.\par\medskip
\begin{lemma}
\label{lemma-base-change-does-not-hold-post}
With $f : X \to S$ and $n$ as in Remark \ref{remark-base-change-holds}
assume $n$ is invertible on $S$ and that for some $q \geq 1$
we have that $BC(f, n, q - 1)$ is true, but $BC(f, n, q)$ is not.
Then there exist a commutative diagram
$$
\xymatrix{
X \ar[d]_f & X' \ar[d] \ar[l] & Y \ar[l]^h \ar[d] \\
S & S' \ar[l] & \Spec(K) \ar[l]
}
$$
with both squares cartesian, where $S'$ is affine, integral, and normal
with algebraically closed function field $K$ and there exists an integer
$d | n$ such that $R^qh_*(\mathbf{Z}/d\mathbf{Z})$ is nonzero.
\end{lemma}

\begin{proof}
First choose a diagram and $\mathcal{F}$ as in
Lemma \ref{lemma-base-change-does-not-hold}.
We may and do assume $S'$ is affine (this is obvious, but
see proof of the lemma in case of doubt).
Let $K'$ be the function field of $S'$
and let $Y' = X' \times_{S'} \Spec(K')$ to get the diagram
$$
\xymatrix{
X \ar[d]_f &
X' \ar[d] \ar[l] &
Y' \ar[l]^{h'} \ar[d] &
Y \ar[l] \ar[d] \\
S &
S' \ar[l] &
\Spec(K') \ar[l] &
\Spec(K) \ar[l]
}
$$
By Lemma \ref{lemma-smooth-base-change-separably-closed}
the total direct image $R(Y \to Y')_*\mathbf{Z}/d\mathbf{Z}$
is isomorphic to $\mathbf{Z}/d\mathbf{Z}$ in $D(Y'_\etale)$; here
we use that $n$ is invertible on $S$.
Thus $Rh'_*\mathbf{Z}/d\mathbf{Z} = Rh_*\mathbf{Z}/d\mathbf{Z}$
by the relative Leray spectral sequence. This finishes the proof.
\end{proof}

\end{document}
